\begin{proof}
Пусть в \(\mathring{B}_R(z_0)\colon \, f(z) = \sum\limits_{n = 0}^{+\infty} c_n (z-z_0)^n\), то
есть главная часть в разложении в ряд Лорана отсутствует. В силу теоремы Абеля, если ряд
сходится в \(\mathring{B}_R(z_0)\), то он сходится и в \(B_R(z_0)\). Его сумма непрерывна в
\(B_R(z_0)\), а значит \(\lim\limits_{z \to z_0} \sum\limits_{n = 0}^{+\infty} c_n (z-z_0)^n =
c_0\), а значит \(z_0\)~--- УОТ.
\end{proof}

\begin{corollary} \label{23:criterion_removable_point}
Из утверждений \ref{23:principal_part_of_bounded_reg_func} и
\ref{23:sufficient_cond_of_removable_point} следует, что:
\begin{enumerate}[wide=20pt, label=\arabic*., itemsep=0ex]
\item \(z_0\)~--- УОТ тогда и только когда главная часть разложения в ряд Лорана
    отсутствует;
\item Если \(f(z)\) ограничена в \(\mathring{B}_R(z_0)\), то \(z_0\)~--- УОТ.
\end{enumerate}
\end{corollary}

\begin{statement}
Если \(z_0 \ne \infty\)~--- полюс, то главная часть ряда Лорана функции \(f(z)\) в \(\mathring{B}
_R(z_0)\) содержит конечное (ненулевое) число слагаемых.
\end{statement}

\begin{proof}
Пусть \(z_0 \ne \infty\)~--- полюс. Тогда \(\exists \mathring{B}_R(z_0) \colon \, \forall z \in
\mathring{B}_R(z_0) \rightarrow |f(z)| < 1\). Обозначим через \(g(z) = \dfrac{1}{f(z)}\). В
\(\mathring{B}_R(z_0)\) имеем \(|g(z)| < 1\), а значит по утверждениям
\ref{23:principal_part_of_bounded_reg_func} и \ref{23:sufficient_cond_of_removable_point} \(z_0\)
~--- УОТ \(g(z)\), а значит главная часть разложения \(g(z)\) отсутствует, то есть \(g(z) =
\sum\limits_{n = 0}^{+\infty} c_n (z-z_0)^n\) в \(\mathring{B}_R(z_0)\). Но по теореме Абеля ряд в
правой части выражения сходится и в полной окрестности \(B_R(z_0)\) к некоторой регулярной функции
\(\hat{g}(z)\colon \, \hat{g}(z) \equiv g(z)\) в \(\mathring{B}_R(z_0)\). Имеем
\(\lim\limits_{z \to z_0} \hat{g}(z) = \lim\limits_{z \to z_0} g(z) = 0\), то есть \(c_0 = 0\).
\(\hat{g}(z) \ne 0\) в \(\mathring{B}_R(z_0)\), то есть \(z = z_0\)~--- нуль конечного порядка
\(n\). Итак \(\exists m \in \N\), такое что в \(B_R(z_0)\) \(\hat{g}(z) = (z-z_0)^m\cdot[c_m +
c_{m+1} (z - z_0) + \dots]\). Ряд в квадратных скобках также сходится в \(B_R(z_0)\) и его сумма~---
регулярная функция \(h(z)\) в \(B_R(z_0)\). \(h(z_0) = c_m \ne 0\), а значит \(w = \dfrac{1}{h(z)}\)
~--- регулярная в \(B_R(z_0)\). \(f(z) = \dfrac{w(z)}{(z-z_0)^m} = \dfrac{a_0 + a_1 (z - z_0) +
\dots}{(z-z_0)^n} = \dfrac{a_0}{(z-z_0)^m} + \dots + a_m + \dots\).
\end{proof}

\begin{definition}
\(z_0\)~--- полюс порядка \(n\), если в разложении в ряд Лорана \ref{eq_23:Laurent_series}
\(c_{-m} \ne 0\) и \(c_n = 0\) для всех \(n < m\).
\end{definition}

\begin{statement}
Если ряд Лорана функции \(f(z)\) в \(\mathring{B}_R(z_0)\) имеет конечное (ненулевое) число
слагаемых в главной части, то \(z_0\)~--- полюс.
\end{statement}

\begin{proof}
\(f(z) = \dfrac{c_{-m}}{(z-z_0)^m} + \dots + \dfrac{c_{-1}}{z-z_0} + c_0 + c_1 (z - z_0) + \dots =
\dfrac{c_{-m} + c_{-m+1}(z-z_0) + \dots}{(z - z_0)^m}\). Степенной ряд числителя сходится и в полной
 окрестности \(B_R(z_0)\). Его сумма непрерывна в этой окрестности, а значит
\(\lim\limits_{z \to z_0} [c_{-m} + c_{-m +1} (z-z_0) + \dots] = c_{-m} \ne 0\). Тогда
\(\lim\limits_{z \to z_0} f(z) = \infty\), то есть \(z_0\)~--- полюс по определению.
\end{proof}

\begin{corollary} \label{23:criterion_pole}
\(z_0 \ne \infty\)~--- полюс \(f(z)\) тогда и только когда главная часть разложения в ряд Лорана
\(f(z)\) в окрестности \(\mathring{B}_R(z_0)\) содержит конечное (ненулевое) число слагаемых.
\end{corollary}

\begin{corollary} \label{23:criterion_essential_point}
Из следствий \ref{23:criterion_removable_point} и \ref{23:criterion_pole} получается, что ИОТ
\(z_0 \ne \infty\) является СОТ тогда и только когда главная часть ряда Лорана имеет бесконечное
число слагаемых.
\end{corollary}

\begin{note}
В силу утверждения \ref{23:ISP_of_1/z} результаты следствий \ref{23:criterion_removable_point},
\ref{23:criterion_pole} и \ref{23:criterion_essential_point} верны и для \(z = \infty\), так как
коэффициенты главной части разложения \(f(z)\) в \(\mathring{B}_R(\infty)\) те же самые, что и
коэффициенты главной части \(f(\dfrac{1}{z})\) в \(\mathring{B}_R(0)\).
\end{note}

\begin{statement}
\textbf{(Теорема Сохоцкого).} Пусть \(z_0 \ne \infty\)~--- СОТ \(f(z)\). Тогда \(\forall z \in \oCo
\; \exists \{z_n\}\colon \, z_n \to z_0, f(z_n) \to A\) при \(n \to \infty\).
\end{statement}

\begin{proof}
\textit{Пусть \(A \ne \infty\).} \textbf{От противного:} пусть такой \(\{z_n\}\) не существует для
этого \(A\). Пусть \(f\)~--- регулярная в \(\mathring{B}_R(z_0)\). Сформулируем некоторое
утверждение:

\begin{equation}
    \forall \delta \in (0, R) \; \forall \epsilon > 0 \; \exists z \in \mathring{B}_\delta(z_0)
    \colon |f(z) - A| < \epsilon. \tag{S} \label{eq_23:auxillary_statement}
\end{equation}

Если \ref{eq_23:auxillary_statement} верно, то существует нужная нам последовательность \(\{z_n\}\)
(достаточно взять \(\delta = \epsilon = \dfrac{1}{n} < R\)). Если такой \(\{z_n\}\) не существует,
то неверно \ref{eq_23:auxillary_statement}, то есть верно его отрицание:

\begin{equation}
    \exists \delta_0 \in (0, R) \; \exists \epsilon_0 > 0 \colon \; \forall z \in
    \mathring{B}_{\delta_0}(z_0) \to |f(z) - A| \ge \epsilon_0. \tag{S'}
    \label{eq_23:negation_auxillary_statement}.
\end{equation}

В \(\mathring{B}_{\delta_0}(z_0)\) определим \(g(z) = \dfrac{1}{f(z) - A}\). В
\(\mathring{B}_{\delta_0}(z_0)\) \(|g(z)| \le \dfrac{1}{\epsilon_0}\), то есть \(g(z)\)~---
ограничена, то есть \(z_0\)~--- УОТ для \(g(z)\), а значит \(\exists \lim\limits_{z \to z_0}
\dfrac{1}{f(z) - A} = B \ne \infty\).

Возможны только два случая:
\begin{enumerate}[wide=20pt, label=\arabic*., itemsep=0ex]
\item \(B \ne 0\). Тогда \(\exists \lim\limits_{z \to z_0} f(z) = A + \dfrac{1}{B} \in \Co\), то
    есть \(z_0\)~--- УОТ для \(f(z)\). Противоречие с условием, что \(z_0\)~--- СОТ для \(f(z)\).
\item \(B = 0\). Тогда \(\exists \lim\limits_{z \to z_0} f(z) + A = \infty\). Так как \(A \ne
    \infty\), то \(\lim\limits_{z \to z_0} f(z) = \infty\), то
    есть \(z_0\)~--- полюс для \(f(z)\). Противоречие с условием, что \(z_0\)~--- СОТ для \(f(z)\).
\end{enumerate}

\textit{Для \(A = \infty\):} если такой \(\{z_n\}\) не существует, то \(f(z)\)~--- ограничена, то
есть\(z_0\)~--- УОТ для \(f(z)\). Получили противоречие.
\end{proof}

\begin{corollary}
В силу утверждения \ref{23:ISP_of_1/z} теорема Сохоцкого остаётся верной и для \(z = \infty\).
\end{corollary}


\Section{Вычеты}

\begin{definition}
Пусть \(z \ne \infty\)~--- ИОТ \(f(z)\). Вычетом функции \(f(z)\) в окрестности \(z_0\) назовём
число \(\res\limits_{z_0} f(z) = c_{-1}\)~--- коэффициент ряда Лорана функции \(f(z)\) в
\(\mathring{B}(z_0)\).
\end{definition}

\begin{statement}
Пусть \(f(z)\)~--- регулярна в \(\mathring{B}_R(z_0)\) и \(\rho \in (0, R)\). Тогда
\(\int\limits_{|z - z_0| = \rho} f(z) \, dz = 2 \pi i \res\limits_{z_0} f(z)\).
\end{statement}

\begin{proof}
В \(\mathring{B}_R(z_0)\) имеем разложение \(f(z) = \sum\limits_{n = - \infty}^{+ \infty} c_n
(z - z_0)^n\). Так как на окружности \(|z - z_0| = \rho\) ряд сходится равноерно, то интегрируя
почленно имеем следующее: \(\int\limits_{|z - z_0| = \rho} f(z) \, dz = \sum\limits_{n = - \infty}^
{+ \infty} c_n \int\limits_{|z - z_0| = \rho} (z - z_0)^n \, dz = c_{-1} \cdot 2 \pi i =
2 \pi i \res\limits_{z_0} f(z)\).
\end{proof}

\begin{note}
Если \(z_0 \ne \infty\) и \(f(z)\) регулярна или \(z_0\)~--- её УОТ, то \(\res\limits{z_0} f(z) =
0\).
\end{note}

\begin{definition}
Пусть \(z_0 = \infty\)~--- ИОТ \(f(z)\). Тогда \(\res\limits_{\infty} f(z) = -c_{-1}\), где
\(c_{-1}\) берётся из разложения \(\sum\limits_{n = - \infty}^{+ \infty} c_n z^n\).
\end{definition}

\begin{note}
Если \(z_0 = \infty\)~--- УОТ, то вполне может быть, что \(\res\limits{\infty} f(z) \ne 0\)
\end{note}

\begin{example}
\(f(z) = \dfrac{1}{z}\). \(z_0 = \infty\)~--- УОТ, \(\res\limits_{\infty} f(z) = -1\).
\end{example}

\begin{statement}
Пусть \(f(z)\) регулярна в \(\mathring{B}_R(\infty)\) и \(R < \rho < + \infty\). Тогда
\begin{equation}
    \int\limits_{|z| = \rho} f(z) \, dz = - 2 \pi i \res\limits_{\infty} f(z).
\end{equation}
\end{statement}

Рассмотрим вычисление вычета в полюсе \(z \ne \infty\).

\begin{statement}
Пусть \(z_0 \ne \infty\)~--- полюс порядка меньше иил равного \(m\) для \(f\). Тогда
\begin{equation}
    \res\limits_{z_0} f(z) = \dfrac{1}{(m-1)!} \lim\limits_{z \to z_0} [(z - z_0)^m f(z)]^{(m - 1)}
\end{equation}
\end{statement}

\begin{proof}
\(f(z) = \dfrac{c_{-m}}{(z - z_0)^m} + \dots + \dfrac{c_{-1}}{z - z_0} + \dots\). При этом здесь
\(c_{-m}\) не обязательно не равно нулю. Тогда \(f(z) (z - z_0)^m = c_{-m} + c_{-m + 1} (z - z_0)
+ \dots + c_{-1} (z - z_0)^{m - 1} + c_0 (z - z_0)^m = \dots\). Ряд в правой части сходится и в
полной окрестности. Его сумма \(g(z)\) регулярна в \(B_R(z_0)\). Тогда
\(f(z)(z - z_0)^m \equiv g(z)\) на \(\mathring{B}_R(z_0)\). Дифференцируя это выражение \((m - 1)\)
раз в \(\mathring{B}_R(z_0)\) имеем, что \([(z - z_0)^m f(z)]^{(m - 1)} \equiv c_{-1} (m-1)! + m!
c_0 (z - z_0) + \dots\). Тогда \(\lim\limits_{z \to z_0} [(z - z_0)^m f(z)]^{(m - 1)} = c_{-1}
(m-1)!\)
\end{proof}

\begin{note}
Для \(n = 1\) формула упрощается до:
\begin{equation}
    \res\limits_{z_0} f(z) = \lim\limits_{z \to z_0} (z - z_0) f(z). \label{eq_24:residue_for_p1p}
\end{equation}
\end{note}

\begin{statement}
Пусть \(f(z) = \dfrac{g(z)}{h(z)}\), где \(g, h\) регулярные в \(z_0\) и \(h\) имеет в \(z_0\)
нуль первого порядка. Тогда \(\res\limits_{z_0} f(z) = \dfrac{g(z_0)}{h'(z_0)}\).
\end{statement}

\begin{proof}
Из \eqref{eq_24:residue_for_p1p} \(\res\limits_{z_0} f(z) = \lim\limits_{z \to z_0} (z - z_0)
\dfrac{g(z)}{h(z)} = \lim\limits_{z \to z_0} \dfrac{g}{\frac{h(z) - h(z_0)}{z - z_0}} =
\dfrac{g(z_0)}{h'(z_0)}\).
\end{proof}

\begin{statement}
\textbf{(Теорема Коши о вычетах).} Пусть \(G\)~--- ОСПГ и \(f(z)\) регулярна в \(\overline{G}\) за
исключением конечного числа ИОТ \(z_1, \dots, z_n \in G\). Тогда:
\begin{equation}
    \int\limits_{\partial G} f(z) \, dz = 2 \pi i \sum\limits_{k = 1}^n \res\limits_{z_n} f(z)
    \label{24_eq:th_koshi_residue}
\end{equation}
\end{statement}

\begin{proof}
Рассмотрим несколько случаев:
\begin{enumerate}[wide=20pt, label=\roman*., itemsep=0ex]
\item Пусть \(G\)~--- ограниченная ОСПГ.

\noindent
\begin{minipage}[c]{0.4\textwidth}
\centering
\begin{tikzpicture}[scale=0.9, >=latex]
    % Внешняя граница области G
    \draw[thick] plot [smooth cycle, tension=0.7] coordinates {
        (0.5, 0.5) (1, 3.5) (3, 4.2) (5.5, 3.8)
        (5.8, 1.2) (4.5, -0.3) (2, -0.6) (0.2, -0.2)
    };
    % Обозначение области
    \node at (6.2, 1.5) {\Large $G$};

    % Заштрихованная область 1 (справа)
    \draw[thick, pattern=vertical lines, pattern color=black!70]
        plot [smooth cycle, tension=0.6] coordinates {
            (4.2, 3) (4.6, 2.5) (4.8, 1.8) (4.2, 1.1) (3.9, 1.8) (4.1, 2.3)
        };

    % Заштрихованная область 2 (слева внизу)
    \draw[thick, pattern=vertical lines, pattern color=black!70]
        plot [smooth cycle, tension=0.7] coordinates {
            (1.8, 1.3) (2.5, 1.3) (2.5, 0.6) (1.8, 0.7)
        };

    % Окрестность точки z_1 (сверху)
    \def\zOne{(2.6, 3.2)}
    \def\rOne{0.75}
    \draw[thick] \zOne circle (\rOne);
    \fill \zOne circle (1.5pt) node[right=1pt] {$z_1$};
    \draw[->] \zOne -- ++(140:\rOne) node[midway, above] {\scriptsize $R_1$};

    % Окрестность точки z_2 (снизу справа)
    \def\zTwo{(3.5, 0.2)}
    \def\rTwo{0.5}
    \draw[thick] \zTwo circle (\rTwo);
    \fill \zTwo circle (1.5pt) node[right=1pt] {$z_2$};
    \draw[->] \zTwo -- ++(120:\rTwo) node[midway, left] {\scriptsize $R_2$};

    % Окрестность точки z_n (снизу слева)
    \def\zN{(1.0, -0.2)}
    \def\rN{0.2}
    \draw[thick] \zN circle (\rN);
    \fill \zN circle (1.5pt) node[right=1pt] {$z_n$};
    \draw[->] \zN -- ++(130:\rN) node[midway, left] {\scriptsize $R_n$};

\end{tikzpicture}
\end{minipage}%
\hfill
\begin{minipage}[c]{0.52\textwidth}
    Выберем \(R_1, \dots, R_n\) так, чтобы \(\overline{B_{R_1}(z_1)}, \dots,
    \overline{B_{R_1}(z_1)} \subset G\) и попарно не пересекались. Рассмотрим область \(D = G
    \setminus (\bigcup\limits_{k = 1}^n \overline{B_{R_k}(z_k)})\). По интегральной теореме Коши для
    ограниченной ОСПГ (\ref{17:int_koshi_ogr_ospg}) \(0 = \int\limits_{\partial D} f(z) \, dz =
    \int\limits_{\partial G} f(z) \, dz = \sum\limits_{k = 1}^n \int\limits_{|z - z_k| = R_k} f(z)
    \, dz = 2 \pi i \sum\limits_{k = 1}^n \res\limits_{z_0} f(z)\). Тогда выполняется
    \eqref{24_eq:th_koshi_residue}.
\end{minipage}
\end{enumerate}
\end{proof}
